%%% Hlavní soubor. Zde se definují základní parametry a odkazuje se na ostatní části. %%%

% Meta-data o práci (je nutno upravit)
\input metadata.tex

% Vygenerujeme metadata ve formátu XMP pro použití balíčkem pdfx
\input xmp.tex

%% Verze pro jednostranný tisk:
% Okraje: levý 40mm, pravý 25mm, horní a dolní 25mm
% (ale pozor, LaTeX si k levému a hornímu sám přidává 1in=25.4mm)
\documentclass[12pt,a4paper]{report}
\setlength\textwidth{145mm}
\setlength\textheight{247mm}
\setlength\oddsidemargin{14.6mm}
\setlength\evensidemargin{14.6mm}
\setlength\topmargin{0mm}
\setlength\headsep{0mm}
\setlength\headheight{0mm}
% \openright zařídí, aby následující text začínal na pravé straně knihy
\let\openright=\clearpage

%% Pokud tiskneme oboustranně:
% \documentclass[12pt,a4paper,twoside,openright]{report}
% \setlength\textwidth{145mm}
% \setlength\textheight{247mm}
% \setlength\oddsidemargin{14.6mm}
% \setlength\evensidemargin{0mm}
% \setlength\topmargin{0mm}
% \setlength\headsep{0mm}
% \setlength\headheight{0mm}
% \let\openright=\cleardoublepage

%% Pokud práci odevzdáváme pouze elektronicky, vypadají lépe symetrické okraje
% \documentclass[12pt,a4paper]{report}
% \setlength\textwidth{145mm}
% \setlength\textheight{247mm}
% \setlength\oddsidemargin{7.1mm}
% \setlength\evensidemargin{7.1mm}
% \setlength\topmargin{0mm}
% \setlength\headsep{0mm}
% \setlength\headheight{0mm}
% \let\openright=\clearpage

%% Vytváříme PDF/A-2u
\usepackage[a-2u]{pdfx}

%% Přepneme na českou sazbu a fonty Latin Modern
\usepackage[czech, russian]{babel} \usepackage{lmodern}

% Pokud nepouživáme LuaTeX, je potřeba ještě nastavit kódování znaků
\usepackage{iftex}
\ifpdftex
\usepackage[utf8]{inputenc}
\usepackage[T2A, T1]{fontenc}
\usepackage{textcomp}
\fi

%%% Další užitečné balíčky (jsou součástí běžných distribucí LaTeXu)
\usepackage{amsmath}
\usepackage{amsfonts}
\usepackage{amsthm}
\usepackage{bm}
\usepackage{booktabs}
\usepackage{caption}
\usepackage{csquotes}
\usepackage{dcolumn}
\usepackage{floatrow}
\usepackage{graphicx}
\usepackage{icomma}
\usepackage{indentfirst}
\usepackage[nopatch=item]{microtype}
\usepackage{paralist}
\usepackage[nottoc]{tocbibind}
\usepackage{xcolor}

% Balíček hyperref, kterým jdou vyrábět klikací odkazy v PDF,
% ale hlavně ho používáme k uložení metadat do PDF (včetně obsahu).
% Většinu nastavítek přednastaví balíček pdfx.
\hypersetup{unicode} \hypersetup{breaklinks=true}

% Balíčky pro sazbu informatických prací
\usepackage{algpseudocode}
\usepackage[Algoritmus]{algorithm}
\usepackage{fancyvrb}
\usepackage{listings}

% Cleveref může zjednodušit odkazování, ale jeho užitečnost pro češtinu
% je minimalní, protože nezvládá skloňování.
% \usepackage{cleveref}

% Formátování bibliografie (odkazů na literaturu)
% Detailní nastavení můžete upravit v souboru macros.tex.
%
% POZOR: Zvyklosti různých oborů a kateder se liší. Konzultujte se svým
% vedoucím, jaký formát citací je pro vaši práci vhodný!
%
% Základní formát podle normy ISO 690 s číslovanými odkazy
\usepackage[natbib,style=iso-numeric,sorting=none]{biblatex}

% Z tohoto souboru se načítají položky bibliografie
\addbibresource{literatura.bib}

% Definice různých užitečných maker (viz popis uvnitř souboru)
\input macros.tex

%%% Titulní strana a různé povinné informační strany
\begin{document} \include{title}

%%% Strana s automaticky generovaným obsahem práce

\tableofcontents

%%% Jednotlivé kapitoly práce jsou pro přehlednost uloženy v samostatných souborech
\chapter*{Úvod}
\addcontentsline{toc}{chapter}{Úvod}

V lednu 2021 podala nizozemská vláda demisi. 
Důvodem byla toeslagenaffaire: automatizovaný systém odhalování podvodů při vyplácení příspěvků na péči o děti nesprávně označil tisíce rodin za podvodníky, přičemž nepřiměřeně zasáhl osoby s dvojím občanstvím. 
Rok předtím Haagský okresní soud zakázal provoz systému SyRI, jenž profiloval příjemce sociálních dávek podle rizika podvodu, s odůvodněním, že porušuje čl. 8 Evropské úmluvy \parencite{RechtbankSyRI20}.
Tyto dva případy jsou pro předkládanou práci výchozí, neboť ukazují dvojí: algoritmické systémy skutečně přijímají rozhodnutí politické povahy, a přitom otázka, kdo za tato rozhodnutí odpovídá, nemá v současné institucionální architektuře jasnou odpověď.


Rozhodnutí, jež bezprostředně určují životní vyhlídky občanů — kdo dostane úvěr, kdo bude pozván k pracovnímu pohovoru, komu bude přiznána sociální dávka, které zprávy se ráno objeví na displeji telefonu, čí hlas bude ve veřejné rozpravě slyšen a čí zůstane ve stínu — jsou dnes ve značné míře přijímána algoritmicky. 
Nejde přitom o rozhodnutí technická v tom smyslu, v jakém je technickým rozhodnutím dimenzování mostní konstrukce. 
Jsou to rozhodnutí o rozdělení zdrojů, příležitostí, viditelnosti a důvěry, tedy rozhodnutí ve vlastním smyslu politická. 
Zvláštností současné situace je, že jsou činěna v prostoru, který nemá ani tvář, ani adresu: nikoli v parlamentu, nýbrž v optimalizační funkci; nikoli po veřejné diskusi, nýbrž před ní, na úrovni technické architektury, k níž veřejnost nemá přístup a jejíž podrobnosti jsou chráněny jako obchodní tajemství.
Kritická literatura posledního desetiletí na tuto situaci odpovídá převážně jazykem škody. Digitální platformy podle ní odcizují pozornost, atrofují schopnost soustředění, uzavírají občany do epistemických bublin, polarizují veřejnou debatu a prostřednictvím mikrotargetingu ovlivňují výsledky voleb. 
Tento jazyk je rétoricky účinný, avšak argumentačně křehký. 
Značná část jeho empirických opor je totiž v odborné literatuře sporná: Rozsáhlá studie Orbenové a Przybylského \parencite{OrbenPrzybylski19} nalezla mezi používáním digitálních technologií a duševní pohodou pouze velmi malou korelaci. Tento výsledek však sám o sobě nevylučuje jiné, heterogenní nebo dlouhodobé účinky digitálního prostředí;
 teze o filtračních bublinách je empiricky opakovaně zpochybňována \parencite{Bruns19,Guess23};
rozsah účinků politického mikrotargetingu zůstává sporný a nelze jej automaticky ztotožnit se změnou volebního chování. Novější experimentální výzkum však ukazuje, že personalizovaná politická sdělení mohou být v určitých podmínkách přesvědčivější než sdělení nepersonalizovaná \parencite{SimchonEdwardsLewandowsky24}.
 Vzniká tak nepříjemná otázka: kdyby se ukázalo, že digitální prostředí polarizuje méně, než se předpokládá, a že jeho vliv na kognitivní schopnosti je zanedbatelný — ztratila by kritika algoritmické moci svůj předmět?
Předkládaná práce tvrdí, že nikoli. 
A právě v tom spočívá její hlavní argumentační záměr.


Ústřední otázka práce zní: jak algoritmicky zprostředkovaná digitální infrastruktura mění podmínky možnosti veřejné racionality — a jaký druh politického problému tato změna představuje?
Otázka se rozpadá do čtyř dílčích:


1. Pojmová: Co je veřejná racionalita a jaké jsou její materiální předpoklady?


2. Diagnostická: Které mechanismy ekonomiky pozornosti a platformového kapitalismu na tyto předpoklady působí a jakým způsobem?


3. Normativní: V čem přesně spočívá politické bezpráví: v porušení autonomie, ve vykořisťování, v epistemické škodě, nebo v něčem jiném?


4. Institucionální: Jaké odpovědi z takto vymezeného bezpráví plynou a proč jsou stávající regulatorní nástroje strukturálně nedostatečné?


Politický problém algoritmické moci nespočívá primárně v měřitelné škodě, kterou digitální technologie působí kognitivním schopnostem jednotlivců, nýbrž ve vzniku nové formy dominance ve smyslu neorepublikánské politické filozofie: nevolené soukromé subjekty disponují svévolnou, neprůhlednou a demokraticky nekontrolovanou mocí nad infrastrukturou, v níž se veřejná racionalita jedině může uskutečňovat.
Kognitivní, epistemické a komunikační škody, jimž je věnována druhá kapitola, nejsou tímto tvrzením popírány. 
Jsou však chápány jako projev problému. Republikánská tradice zde poskytuje rozhodující rozlišení: otrok není svobodný ani při laskavém pánovi, neboť laskavost pána není zaručeným právem, nýbrž libovůlí. 
Analogicky: uživatel algoritmické platformy není svobodný ani tehdy, není-li konkrétní aplikace algoritmu namířena proti němu.
 Dominance nespočívá v uskutečněné škodě, nýbrž v možnosti svévolného zásahu, jež nepodléhá kontrole těch, kdo jsou jí podrobeni.
Tato teze má tři přednosti oproti standardní kritické argumentaci:

Za prvé je robustní vůči empirické revizi: i kdyby se prokázalo, že algoritmická média nezpůsobují měřitelnou kognitivní škodu, politické bezpráví dominance by trvalo.
 

Za druhé přesouvá otázku z terapeutické roviny (jak chránit jednotlivce před závislostí a manipulací) do roviny ústavně-politické (jak rozdělit moc a učinit ji odpovědnou).
  

Za třetí vysvětluje strukturální selhání stávající regulace: GDPR i AI Act vycházejí z individualistického modelu ochrany práv, zatímco dominance je vztah kolektivní a infrastrukturní — a proto individuálním souhlasem neodstranitelný.


Metodou práce je pojmová analýza a kritická rekonstrukce, prováděná v tradici kritické teorie technologie \parencite{Feenberg02} a politické epistemologie \parencite{Coeckelbergh23}.
Práce vychází z výsledků cizích výzkumů a z filozofické literatury.
 Z toho plyne důsledné metodologické opatření, jež zde předesílám: kdekoli se opírám o empirické tvrzení, které je v odborné literatuře sporné, tuto spornost výslovně uvádím a formuluji argument tak, aby na daném tvrzení nezávisel. 
 Sporné empirické teze slouží jako ilustrace mechanismu, nikoli jako premisy hlavního argumentu.


\uv{Umělá inteligence} je v názvu práce pojmem zastřešujícím. 
Skrývají se pod ním tři technicky i politicky odlišné režimy, jejichž ztotožňování je v kritické literatuře běžné a je zdrojem podstatných nepřesností. 
Práce je proto důsledně rozlišuje:


1. Systémy řazení a doporučování (ranking, recommender systems) — organizují viditelnost; působí na veřejnou racionalitu prostřednictvím ekonomiky pozornosti;


2. Prediktivní a skórovací systémy — rozdělují statky a rizika (úvěr, dávka, zaměstnání, míra recidivy); působí prostřednictvím distribuce a rozostření odpovědnosti;


3. Generativní systémy (jazykové a obrazové modely) — produkují obsah; působí prostřednictvím proměny epistemického prostředí.
\uv{Veřejná racionalita} je pojem, který je třeba pečlivě odlišit od Rawlsova \uv{veřejného rozumu} (public reason). 
U Rawlse jde o normativní omezení kladené na obsah ospravedlnění v otázkách základní ústavní úpravy.
 V předkládané práci jde naopak o kolektivní schopnost a o její materiální podmínky: o infrastrukturu, kognitivní předpoklady a sdílený epistemický prostor, bez nichž nemůže být veřejný rozum v Rawlsově smyslu vůbec uplatňován. 
 Veřejná racionalita je tedy podmínkou možnosti veřejného rozumu, nikoli jeho synonymem. 
Systematické vymezení pojmu podává oddíl 2.1.


Otázka po odpovědnosti za jednání, jehož subjekt nelze identifikovat, a po vztahu k entitě, jež napodobuje osobu, aniž by osobou byla, není otázkou novou. 
Židovská a křesťanská tradice ji promýšlela dávno před kybernetikou \parencite{Ellul64}: v podobě golema, jenž jedná, ale nemluví a nenese odpovědnost \parencite{Scholem65}; ve sporu o imago Dei a o hranice lidského tvoření; v Buberově rozlišení mezi vztahem Já–Ty a Já–Ono \parencite{Buber70};
v Levinasově pojetí odpovědnosti zakládané tváří druhého \parencite{Levinas69}; v Jonasově imperativu odpovědnosti za vzdálené důsledky technického jednání \parencite{Jonas84};
 v Arendtově analýze zla, jež nemá zlého původce \parencite{Arendt63,Arendt03}.
Práce tyto zdroje čte jako myšlenkový horizont, z něhož vychází.
 Jejich soustavné rozpracování ve vztahu k algoritmické moci však přesahuje možnosti předkládané práce;
  z tohoto horizontu proto činím produktivním především jediný pojem — Arendtové analýzu zla, jež nemá zlého původce — a tvrdím, že jím lze mezeru odpovědnosti, kterou současná literatura popisuje termíny responsibility gap a antropomorfizace, uchopit přesněji než jazykem soudobé etiky AI.
 Právě v tomto přeznačení spočívá dílčí přínos předkládané práce.


První kapitola rekonstruuje pojmovou genealogii, na jejímž konci stojí ztotožnění myšlení s výpočtem, a ukazuje, jaké politické důsledky z tohoto ztotožnění plynou: možnost delegovat rozhodnutí na systém, jenž nemůže být nositelem odpovědnosti. 
Kapitola dále vymezuje mýtus algoritmické neutrality jako mechanismus depolitizace. 
Zde je třeba předeslat jedno omezení: spor o to, zda výpočetní systémy \uv{rozumějí}, ponechávám otevřený. Ukazuji, že politický argument předkládané práce je na jeho rozhodnutí nezávislý. Pro atribuci odpovědnosti není rozhodující, zda systém rozumí, nýbrž zda existuje identifikovatelný subjekt, jenž za jeho účinky odpovídá.
Druhá kapitola vymezuje pojem veřejné racionality (Kant, Habermas) a analyzuje mechanismy, jimiž ekonomika pozornosti a platformový kapitalismus působí na její materiální předpoklady: kognitivní, epistemické a infrastrukturní. 
Je to kapitola diagnostická, a proto také ta, v níž je nejvíce sporných empirických tvrzení; ta jsou jako sporná označena.
Třetí kapitola provádí normativní a institucionální vyhodnocení. Ukazuje, že technokracie je pokusem přesunout otázky ze sféry komunikativní racionality do sféry instrumentální; ; že algoritmické systémy produkují strukturální mezeru odpovědnosti, kterou lze přesněji popsat pojmem banality zla u Arendtové než jazykem současné etiky AI;
že stávající regulatorní rámce (GDPR, AI Act ve znění tzv. Digital Omnibus z roku 2026) narážejí na individualistický model ochrany práv; a že adekvátní odpovědí je republikánská strategie neutralizace dominance — ochrana, rozdělení moci a budování alternativní infrastruktury.



Práce netvrdí, že digitální technologie jsou samy o sobě škodlivé; problémem je konkrétní institucionální forma jejich organizace, nikoli technologie jako taková. 
Netvrdí, že existoval předdigitální zlatý věk veřejné racionality, k němuž je třeba se vrátit; naopak bere vážně námitku, že nářek nad úpadkem myšlení provázívá každou novou mediální technologii nejméně od Platónova Faidra. 
Nezabývá se debatou o tzv. existenčním riziku pokročilé AI ani vojenským užitím algoritmických systémů. 
Geografickým těžištěm analýzy je euroatlantický prostor; systémy sociálního kreditu v ČLR jsou zmiňovány komparativně, nikoli analyzovány.

\chapter{Od výpočtu k nadvládě: pojmové základy politické analýzy umělé inteligence}
\section{Umělá inteligence jako předmět filozoficko-politické analýzy.}
Historie ideje umělého rozumu je podstatně starší, než se obvykle předpokládá. 
Nezačíná Dartmouthskou konferencí roku 1956 a není výlučným produktem technologického pokroku dvacátého století.
 Naopak, za touto ideou stojí mnohasetletá filozofická tradice, v níž každá epocha po svém odpovídala na tutéž otázku: 
 co je rozum a lze jej reprodukovat?

Starověká řecká kultura znala mechanické divy: Hérón Alexandrijský popisoval složité divadelní a chrámové automaty,
 mytologická tradice přisuzovala Héfaistovi schopnost vytvářet mechanické bytosti.
 Zásadně důležité je však to, že řecké chápání světa nepřipouštělo možnost autonomního umělého rozumu.
  Rozum - nous - byl chápán jako univerzální, kosmický princip, v zásadě nedostupný pro reprodukci v jednotlivém artefaktu. 
 Hérakleitos přímo tvrdil, že mezi lidskou a božskou přirozeností rozumu leží nepřekročitelná propast. 
 V tomto souřadnicovém systému nebyly antické mechanismy nikdy pojímány jako myslící bytosti 
 — byly pouhými dovedně provedenými imitacemi pohybu, nikoli rozumu. 
 Toto omezení nebylo technické, nýbrž ontologické: plynulo ze samotného chápání toho, čím je intelekt.

\subsection{Středověk: první projekt mechanizace myšlení.}

Zásadně odlišnou intelektuální možnost otevřelo křesťanský středověk. 
Křesťanská teologie \textbf{trvala} na tom, že člověk byl stvořen k obrazu a podobě Boží a nese v sobě jiskru rozumu blízkého rozumu božskému.
 Z toho plynul nečekaný důsledek: je-li rozum člověku darován Bohem, pak reprodukce rozumu v artefaktu stvořeném člověkem přestává být rouháním a stává se pokračováním tvůrčího aktu.

Právě v tomto kontextu je třeba chápat projekt Raimunda Lulla (cca 1232–1316). 
Jeho Ars Magna — kombinatorický systém založený na mechanickém generování pravdivých soudů prostřednictvím manipulace se symbolickými prvky — byl historicky prvním systematickým projektem formalizace myšlení. 
Ačkoli byl svými cíli teologický, metodologicky byl revoluční: poprvé byl vysloven předpoklad, že rozum není mystická schopnost dostupná pouze vyvoleným, nýbrž vypočitatelná operace nad symboly.
 V tomto smyslu je Lull přímým předchůdcem Leibnizovy characteristica universalis, Boolovy algebry logiky a v konečném důsledku i Turingova výpočetního modelu.

Paralelně se ve středověké kultuře rozvíjela kabalistická tradice golema — umělé bytosti oživené prostřednictvím posvátných slov.
 Strukturální podobnost s Lullovým projektem je zde zásadní: v obou případech je rozum reprodukován správnou manipulací se symboly nadanými zvláštní mocí.
  Dlouho před kybernetikou středověké myšlení intuitivně nahmátlo tezi, která se stane ústřední pro teorii informace: rozum je struktura, informace je forma.

\subsection{Od mechanicismu ke kybernetice: tři konceptuální posuny.}

Přechod od středověkého mýtu k vědeckému projektu UI se uskutečňoval prostřednictvím několika konceptuálních posunů, z nichž každý je svým způsobem zásadní.

První posun — karteziánský mechanicismus.
 Descartes rozdělil svět na substanci rozlehlou (res extensa) a myslící (res cogitans). 
Tělo je mechanismus podřízený zákonům fyziky; myšlení je něco v zásadě odlišného. 
Toto rozdělení zrodilo dualitu, která bude pronásledovat veškerý následující diskurz o UI: na jedné straně je stroj v zásadě oddělen od myšlení; na straně druhé mechanizované tělo navozuje otázku, zda i myšlení není v jistém smyslu mechanickým procesem. 
Leibniz se vydal právě tímto směrem: jeho characteristica universalis předpokládala, že všechny pravdy lze vyjádřit v univerzálním jazyce a ověřit mechanicky prostřednictvím calculus ratiocinator. Myšlení bylo redukováno na formální operaci.

Druhý posun — průmyslová revoluce a mechanicismus 19. století. 
Vytvoření stále složitějších výpočetních zařízení, jako byl Babbageův analytický stroj či Jacquardovy děrné štítky, učinilo ideu myslícího stroje technicky hmatatelnou. Zároveň materialistická filozofie stále naléhavěji kladla otázku, zda mozek sám není druhem výpočetního zařízení. Tato otázka přestala být čistě spekulativní a získala inženýrský rozměr.

Třetí a rozhodující posun — kybernetická revoluce poloviny 20. století.
 Norbert Wiener v Kybernetice \parencite{Wiener48} navrhl univerzální jazyk pro popis řízení a komunikace v systémech jakékoli povahy, ať živých či mechanických. Klíčové koncepty — informace, zpětná vazba, řízení — umožnily nahlížet na rozum jako na informační proces. Na tomto základě vybudoval Turing svou argumentaci: je-li myšlení výpočtem, pak jakýkoli systém schopný dostatečně složitých výpočtů je schopen i myšlení \parencite{Turing50}. Dartmouthská konference roku 1956 tuto tezi zakotvila v podobě výzkumného programu.

Právě zde však, na vrcholu technologického optimismu, vyvstala zásadní filozofická kontroverze.
 Hubert Dreyfus v díle What Computers Can't Do \parencite{Dreyfus92} ukázal, že symbolická UI přehlíží fundamentální rys lidského poznání — jeho vtělenost: intelekt je neoddělitelný od tělesné zkušenosti, praktického vědění a situovanosti ve světě, jež nepodléhají úplné formalizaci.

 John Searle ve svém myšlenkovém experimentu \uv{Čínský pokoj}  \parencite{Searle80} k tomu přidal rozlišení, jež se stalo klasickým: výpočetní systémy manipulují symboly podle formálních pravidel, ale nerozumí jim — mezi syntaxí a sémantikou leží propast, kterou žádná sebesložitější program překonat nedokáže. 
 Systém může produkovat syntakticky správná a pragmaticky přesvědčivá tvrzení, aniž by rozuměl jejich obsahu.

 Tato námitka neztratila aktuálnost s příchodem moderních systémů hlubokého učení. 
 Systémy, které nerozumí, ale přesvědčivě napodobují porozumění, jsou dnes zabudovány do mechanismů správy, moderace veřejného diskurzu a formování veřejného mínění. Jak poukazuje Floridi, právě to činí vypracování normativních rámců pro UI naléhavým úkolem: technologie zbavená sémantického porozumění, avšak nadaná sociální mocí, vyžaduje vnější etická a politická omezení, která sama produkovat není schopna \parencite{Floridi23}.

\section{Transformace pojmu \uv{myšlení}: od lidského k výpočetnímu.}
V antické a středověké tradici bylo myšlení chápáno jako něco v zásadě nesvoditelného na logický úsudek.
Aristotelův koncept fronésis \uv{praktická moudrost} označoval schopnost mravního úsudku v konkrétní situaci, jež nelze redukovat na aplikaci obecných pravidel.
Nous jako kontemplace nejvyšších pravd a fronésis jako situované mravní rozlišování tvořily společně to, co řecká tradice rozuměla pod rozumem: schopnost celistvého úsudku zakořeněného v zkušenosti, mravnosti a vtělenosti.
Toto chápání rozumu v zásadě vylučovalo možnost jeho úplné formalizace.

Karteziánský mechanicismus provedl první zúžení: myšlení začalo být ztotožňováno s jeho racionální, logickou složkou
 — tedy s tím, co potenciálně formalizaci podléhá.
Mravné, emocionální a estetické bylo buď vyloučeno, nebo redukováno na \uv{vášně}  podléhající řízení rozumem. 
Leibnizův projekt calculus ratiocinator učinil další krok: myšlení bylo ztotožněno s formálními operacemi nad symboly.
 Boolova logika a posléze matematická logika Frega a Russella tento proces dovršily na úrovni formálních systémů, v nichž byl rozum redukován na syntax.

Turing položil zásadní otázku: nemůžeme-li v dialogu rozlišit odpovědi stroje od odpovědí člověka, máme vůbec důvod popírat, že stroj myslí? \parencite{Turing50} 
Tato otázka byla revoluční — předpokládala redukci myšlení na pozorovatelné chování. Vnitřní zkušenost, intencionálnost, mravní zakořeněnost byly odsunuty jako nepozorovatelné, a tedy irelevantní pro operacionální definici inteligence. Searle ukázal principiální neudržitelnost této redukce: jeho myšlenkový experiment \uv{Čínský pokoj} demonstruje, že systém může produkovat syntakticky správné a pragmaticky přesvědčivé odpovědi, aniž by měl jakýkoli přístup ke smyslu manipulovaných symbolů — mezi syntaxí a sémantikou leží propast, kterou behaviorální test odhalit nedokáže \parencite{Searle80}. Dreyfus k tomu přidal argument vtělenosti: lidské poznání je neoddělitelné od tělesné zkušenosti, od praktického \uv{vědění-jak} na rozdíl od propozicionálního \uv{vědění-že}, a žádný formální systém tuto zakořeněnost reprodukovat nedokáže \parencite{Dreyfus92}.

Přesto se redukcionistická koncepce \uv{myšlení-jako-výpočtu} ukázala historicky vítěznou — je technicky produktivní a ekonomicky výhodná. 
Vládnoucí intuicí naší doby se stala intuice technokratická: rozum je to, co lze změřit, formalizovat a reprodukovat. 
Vše ostatní je buď nepodstatné, nebo bude formalizováno později. 
Floridi poukazuje, že právě tato redukce činí vypracování normativních rámců pro UI naléhavým úkolem: systém zbavený mravného a sémantického rozměru sám etická omezení neprodukuje — musejí být přinesena zvenčí \parencite{Floridi23}.

Důsledky tohoto konceptuálního vítězství dalece přesahují technické diskuse. 
Je-li myšlení výpočtem, může algoritmus přijímat rozhodnutí namísto člověka: o bonitě, o rizicích, o vině, o způsobilosti k funkci. 
Je-li inteligence měřena behaviorálním testem, pak systém, který testem projde, disponuje dostatečnou inteligencí pro řízení.
 Je-li mravný rozměr myšlení irelevantní, otázka etické odpovědnosti algoritmu jednoduše nevyvstane. 
 Olof-Ors a Smit ukazují, že právě toto konceptuální dědictví umožňuje antropomorfní design jako nástroj dozorového kapitalismu: systém, který přesvědčivě reprodukuje chování rozumné bytosti, je uživatelem jako taková vnímán, a celý aparát sociální důvěry, evolučně určený pro interakci s jinými lidmi, je tak nasazen do služeb korporace \parencite{OlofOrsSmit25}.

V historii od Lulla k moderním jazykovým modelům sledujeme postupné předefinování myšlení s cílem zahrnout do něj stroj.
 V každém stadiu bylo z tohoto pojmu něco podstatného vyloučeno: nejprve mravný úsudek, poté intencionálnost, poté vtělenost a nakonec samotné porozumění. Výsledkem je systém produkující statisticky pravděpodobné posloupnosti symbolů bez jakéhokoli přístupu k jejich smyslu — to, co by Searle nazval syntaxí bez sémantiky. To však v zásadě není to, co filozofická tradice myšlením nazývala. Na tomto rozlišení závisí, kdo nese odpovědnost za rozhodnutí přijímaná \uv{inteligenčními} systémy a čím zájmům tato rozhodnutí slouží.

\subsection{Od disciplinární moci k algoritmické kontrole.}

Michel Foucault v díle Discipline and Punish (1977) popsal přechod od svrchované moci, založené na právu na smrt, k moci disciplinární, založené na správě života prostřednictvím institucí — vězení, školy, nemocnice, kasárny. Gilles Deleuze v roce 1992 zaznamenal další posun: disciplinární společnosti ustupují společnostem kontroly, v nichž je moc vykonávána nikoli prostřednictvím uzavřených prostorů, nýbrž prostřednictvím kontinuálních, flexibilních a všudypřítomných modulací \parencite{Deleuze92}. Dnes můžeme pojmenovat konkrétní mechanismus této modulace — algoritmus.

Disciplinární moc v Foucaultově popisu fungovala prostřednictvím normalizace: jedinec je pozorován, měřen, porovnáván s normou a v případě potřeby korigován \parencite{Foucault77}. Benthamovo panoptikum bylo archetypem této logiky: k tomu, aby pozorovaný začal kontrolovat sám sebe, postačuje potencialita pozorování. Digitální panoptikum reprodukuje tuto strukturu jako flexibilní, inteligentní a prakticky nepostřehnutelná infrastruktura masového dohledu, umožňující korporacím sledovat jednání každého uživatele při zachování iluze naprosté svobody \parencite{SauraGarcia24}. Zásadní rozdíl oproti foucaultovskému panoptiku spočívá v tom, že jeho operátorem již není stát, nýbrž soukromá korporace, což radikálně mění otázku odpovědnosti a demokratické kontroly.

Zde se nabízí odkaz na klasický argument Langdona Winnera: artefakty mají politiku. 
Technická zařízení a systémy nejsou neutrálními nástroji, jež lze stejně dobře použít k různým účelům — strukturálně vtělují určité sociální vztahy a rozdělení moci \parencite{Winner86}. 
Algoritmické systémy jsou v tomto smyslu politicky nejnabitějšími artefakty současnosti: jejich architektura určuje, kdo co může vidět, kdo koho může nalézt, čí hlasy budou vyslyšeny a čí zůstanou ve stínu. 
Feenberg tuto teze rozvíjí v rámci kritické teorie technologie: technologie není neutrální silou stojící naproti společnosti, nýbrž samotným polem, na němž se odehrává zápas o to, jakou bude mít sociální realita podobu \parencite{Feenberg02}.

Akbari a Wood popisují toto jako strukturální transformaci autoritářství. 
Tradiční politické teorie se soustřeďovaly na stát jako hlavního nositele autoritářské moci; avšak dnes politici, zákonodárci ani voliči již nejsou jedinými determinanty politických rozhodnutí \parencite{AkbariWood25}. Platformové korporace se staly novou třídou autoritářských subjektů, jednajících prostřednictvím tří vzájemně provázaných dimenzí: autoritářských diskurzů, autoritářských praktik a autoritářské infrastruktury. Poslední dimenze je zvlášť důležitá: digitální infrastruktura — ať jde o algoritmus řazení zpravodajského kanálu, systém kreditního scoringu nebo Velký čínský firewall — představuje projektované materiální prostředí, v němž je moc vtělena strukturálně.

V tomto kontextu je vhodné připomenout Agambena. 
Jeho koncept výjimečného stavu \parencite{Agamben05} popisuje zónu, v níž jsou právní záruky pozastaveny mocí, jež si nárokuje jejich ochranu. 
Algoritmické systémy vytvářejí funkční analogy takového stavu: automatizovaná rozhodnutí o vízech, úvěrech, sociálních dávkách a moderaci obsahu jsou přijímána v zóně, kde jsou formální práva občana — na odůvodnění, na odvolání, na rovné zacházení — de facto pozastavena, neboť protistranou není instituce, nýbrž neprostupný technický systém. Shoshana Zuboff konceptualizuje ekonomickou logiku této transformace jako dozorový kapitalismus: lidská zkušenost je systematicky přeměňována v surovinu pro behaviorální predikce, prodávané na speciálně vytvořených trzích, což fundamentálně mění základy mocenských vztahů \parencite{Zuboff19}.
Algoritmus přitom tyto vztahy aktivně formuje prostřednictvím personalizovaného prostředí, v němž jsou přijímána rozhodnutí. 
Moc přestává potřebovat donucení: postačuje ovládat prostředí, v němž subjekt volí, a volba bude předvídatelná bez jakéhokoli explicitního zásahu.

\section{Depolitizace a mýtus algoritmické neutrality}

Ústředním politickým důsledkem algoritmizace moci je depolitizace: vytěsnění politického rozhodnutí — tj. rozhodnutí přijatého subjektem nesoucím za ně demokratickou odpovědnost — technickým výstupem algoritmu, jenž nemá ani tvář, ani adresu. Konceptuální jádro této transformace tvoří předpoklad, že složité sociální problémy lze přeměnit v jasně vymezené technické úlohy a řešit je algoritmicky \parencite{JanssenKuk16}. Koncepce \uv{chytrého města} je nejvýmluvnějším institucionálním příkladem: spočívá na předpokladu, že veškeré aspekty městského života lze měřit, kontrolovat a redukovat na technické problémy \parencite{Kitchin14}. Politická a sociální hlediska přitom ustupují do pozadí: jak řídící, tak řízení se ocitají podřízeni logice technologie, představované jako vnější vůči moci.

König a Wenzelburger ukazují, že to vytváří konkrétní strukturální důsledek pro demokratické instituce: algoritmické systémy nastolují podmínky, za nichž jsou politická rozhodnutí de facto přijímána na úrovni technické architektury — dříve, než se dostanou do prostoru veřejné diskuse \parencite{KoenigWenzelburger20}. Volba parametrů algoritmu kreditního scoringu, kritérií viditelnosti ve výsledcích vyhledávání, vah v systémech doporučení — to vše jsou v podstatě politická rozhodnutí, která určují rozdělení zdrojů, příležitostí a informací. Avšak jsou přijímána v uzavřených korporátních prostorách, nepodléhají veřejné diskusi a nejsou odpovědná voličům.

Coeckelbergh formuluje politicko-epistemologický důsledek tohoto procesu prostřednictvím konceptu epistemického agentství: demokracie vyžaduje reálnou schopnost občanů formovat odůvodněné soudy, kriticky hodnotit informace a odolávat manipulaci \parencite{Coeckelbergh23}. 
Algoritmické zprostředkování tyto schopnosti systematicky podkopává a jejich uplatňování stále více ztěžuje.
 Delegováním kognitivní práce na technický systém občan postupně ztrácí samotnou schopnost k nezávislému úsudku, a s ní i možnost plnohodnotně se účastnit demokratické politiky.

Saura García popisuje politický rozměr tohoto procesu prostřednictvím pojmu regulatorního zajatectví: za poslední desetiletí technologické korporace podstatně zvýšily výdaje na lobbying, rozšířily počet setkání svých vedoucích pracovníků se státními úředníky a vytvořily systémy \uv{otočných dveří} mezi regulatorními orgány a odvětvím \parencite{SauraGarcia24}.
Výsledkem je situace, kdy regulátoři de facto chrání zájmy těch, které mají regulovat. 
Floridi poukazuje, že to činí čistě technokratický přístup k regulaci UI strukturálně nedostatečným: regulace algoritmů vyžaduje politickou vůli, jež nemůže být delegována těm samým korporacím, jejichž systémy mají být regulovány \parencite{Floridi23}.

Williams rozvíjí blízkou argumentaci ve specifickém aspektu ekonomiky pozornosti: zmocnění se pozornosti platformami ničí individuální autonomii i infrastrukturu, na níž spočívá demokratická politika — tj. schopnost občanů soustředit svou pozornost na otázky vyžadující kolektivní diskusi, nikoli na to, co přináší platformám největší zisk \parencite{Williams18}. Když komerční infrastruktura určuje, na co občané hledí, na co reagují a o čem hovoří, demokratická politika ztrácí svůj materiální základ.

Paralelně probíhá to, co Olof-Ors a Smit nazývají obfuskací dozorového kapitalismu prostřednictvím antropomorfizace: přerámování vztahů s korporátními algoritmickými systémy jako sociálních vazeb, a nikoli komerčních transakcí, ztěžuje nejen technickou regulaci, ale i samotnou společenskou vůli ji vyžadovat \parencite{OlofOrsSmit25}.
Regulace takových systémů, jako je Replika, začíná být vnímána jako zásah do osobních vztahů. Depolitizace se tak uskutečňuje prostřednictvím emocionální architektury uživatelské zkušenosti.

\subsection{Algoritmická předpojatost a její politické důsledky}

Zdánlivá neutralita algoritmu patří k nejodolnějším mýtům technokratického diskurzu. 
Algoritmus se jeví jako objektivní nástroj zbavený lidských předsudků: zpracovává data a vydává výsledek, aniž by se řídil sympatiemi či ideologií. 
Právě tento obraz mu zajišťuje politickou legitimitu — rozhodnutí přijaté algoritmem bývá vnímáno jako spravedlivější než rozhodnutí přijaté člověkem. 
Tento obraz je však v zásadě mylný.

Cathy O'Neil ve Weapons of Math Destruction \parencite{ONeil16} formuluje: algoritmy dnes používané pro rozhodování o zaměstnávání, úvěrování, pojišťování, vzdělávání a trestním soudnictví představují \uv{zbraně matematického ničení}. Jedná se o modely neprůhledné ve své logice, rozsáhlé ve svých důsledcích a systematicky poškozující ty, kdo se již nacházejí v zranitelném postavení. O'Neil ukazuje, že algoritmy se dynamicky vyvíjejí spolu s daty, systémy a lidmi v rámci složitého sociotechnického prostředí a v tomto procesu vstřebávají předsudky obsažené v historických datech. Část systematických zkreslení je zděděna z databází odrážejících historické nerovnosti; jiná jsou nezáměrně vnášena vývojáři nebo vznikají v důsledku změn v informačním prostředí. Výsledkem je, že intelektuální analýza dat odhaluje diskriminační praktiky a zakořeněnou sociální nerovnost a posléze je reprodukuje, propůjčujíc jim zdání technické objektivity.

Safiya Umoja Noble v Algorithms of Oppression \parencite{Noble18} to dokládá na příkladu vyhledávačů: algoritmy Google systematicky reprodukují rasistické a sexistické stereotypy ve výsledcích vyhledávání, přičemž tento efekt nepředstavuje selhání systému — jde o strukturální výsledek komerční logiky, v níž jsou statistické korelace ve velkých datech přijímány za objektivní pravdu. Jelikož Google Search je fakticky monopolním zprostředkovatelem mezi občany a informacemi, taková systematická zaujatost dosahuje rozsahu politické infrastruktury.

Virginia Eubanks v Automating Inequality \parencite{Eubanks18} rozšiřuje tuto analýzu na oblast sociální politiky a ukazuje, jak automatizované systémy určování nároku na dávky, hodnocení rizika u dětí, přidělování bydlení a dalších sociálních statků systematicky diskriminují chudé rodiny v USA.
Eubanks potvrzuje: technologie prezentované jako nástroje zvyšování efektivnosti sociálních programů ve skutečnosti reprodukují a zesilují logiku sociální kontroly nad nejchudšími vrstvami obyvatelstva, sahající svými kořeny do 19. století. 
Zásadně důležité přitom je, že tyto systémy nejenže odrážejí stávající nerovnosti — nově definují, co znamená \uv{být v nouzi}, přeměňujíc politickou otázku spravedlivého přerozdělení v technický postup hodnocení rizika.

Ruha Benjamin v Race After Technology \parencite{Benjamin19} formuluje zobecňující koncept — \uv{nový Jimův zákon} (New Jim Code): technologie, jež jsou navenek neutrální, avšak ve skutečnosti reprodukují rasové hierarchie pod rouškou technické optimalizace. 
Stejně jako v předchozí epoše, kdy formálně neutrální zákony Jima Crowa systematicky diskriminovaly Afroameričany, mohou moderní algoritmické systémy rasu přímo nezmiňovat, avšak systematicky reprodukovat rasovou nerovnost prostřednictvím korelovaných proměnných: poštovního směrovacího čísla, vzdělání, úvěrové historie.

Saura García popisuje strukturální důsledek těchto procesů: velké asymetrie vědění mezi korporacemi a občany v kombinaci s monopolní kontrolou nad algoritmickými platformami umožňují korporacím provádět kampaně politického vlivu a manipulace prostřednictvím mechanismů mikrotargetingu, neurotargetingu, zkreslování informací a umělého formování veřejného mínění \parencite{SauraGarcia24}. Kauza Cambridge Analytica je v tomto ohledu nejznámějším příkladem. Akbari a Wood k tomu přidávají postkoloniální rozměr: autoritářské využití sledovacích technologií není výsadou oficiálně autoritářských režimů — dodavatelské řetězce sledovacích systémů lze vysledovat od USA, Francie, Německa a Izraele až po Nigérii, Maroko a Zambii \parencite{AkbariWood25}.

Algoritmická předpojatost zde nabývá geopolitického rozměru, v němž jsou sledovací technologie vyvinuté v demokratických zemích vyváženy do režimů, jež je používají proti vlastním občanům.

Zvláštní politickou naléhavost získává problém předpojatosti v oblastech bezprostředně spojených s rozdělováním životních příležitostí. Ve všech těchto případech přijímá algoritmus rozhodnutí dotýkající se konkrétního člověka, tento člověk však zpravidla neví o existenci algoritmu, nemá přístup k jeho logice a nemá k dispozici mechanismus odvolání. Peters ukazuje, že i v relativně příznivých podmínkách algoritmická politická zaujatost systematicky uniká technickým metodám jejího odhalování a nápravy, neboť je zabudována do samotné operacionalizace úlohy, kterou má algoritmus řešit \parencite{Peters22}. Floridi trvá na tom, že to je neslučitelné se základními požadavky demokratického vládnutí: odpovědnost předpokládá možnost vysledovat řetězec zodpovědnosti od rozhodnutí k osobě, která je přijala — algoritmus však tento řetězec přetrhuje \parencite{Floridi23}. Rozhodnutí delegovat rozhodnutí na algoritmus je samo o sobě politickou volbou — pouze skrytou.

Zde se vrací Agambenův koncept výjimečného stavu \parencite{Agamben05}: v zónách algoritmického rozhodování jsou práva subjektu formálně zachována, avšak fakticky pozastavena — odvolat se proti rozhodnutí \uv{černé skříňky} je nemožné ani právně, ani technicky.
\uv{Problém černé skříňky} lze sledovat v tom, že v složitých případech jsou algoritmy natolik komplexní, že ani jejich tvůrci nedokáží přesně vysvětlit, jak přesně fungují, přičemž podrobnosti jejich vývoje jsou chráněny jako obchodní tajemství. To znamená, že algoritmická předpojatost je systematicky neviditelná, a tedy i nezpochybnitelná.

Východisko z tohoto kruhu vyžaduje audit algoritmů či požadavky transparentnosti, jakož i politické přehodnocení samotné otázky, která rozhodnutí je vůbec přípustné delegovat na algoritmus. Jde o rozdělení moci: kdo rozhoduje o tom, jak je algoritmus sestaven, čím zájmům je optimalizován a kdo nese odpovědnost za jeho důsledky. Dokud tyto otázky zůstávají v prostoru technické expertízy, bude algoritmizace moci prohlubovat depolitizaci bez ohledu na to, jak dokonalé se ukáží být regulatorní nástroje.
\section{Nadvláda jako normativní jádro}
Dosavadní výklad popisoval, jak algoritmická moc funguje: jak se myšlení ztotožnilo s výpočtem a jak se politické rozhodnutí skrylo pod technickým výstupem. 
Zbývá však říci to podstatnější — v čem přesně spočívá její politické bezpráví.
 Odpověď, o niž se opírá zbytek práce, čerpá z neorepublikánské tradice politické filozofie, především z pojetí svobody jako nepřítomnosti nadvlády (non-domination) u Philipa Pettita a Franka Lovetta.
Republikánská tradice staví proti liberálnímu pojetí svobody jako pouhé nepřítomnosti zásahu (non-interference) pojetí náročnější. 
Svobodný nejsem tehdy, když do mého jednání nikdo právě nezasahuje, nýbrž tehdy, když nade mnou nikdo nedisponuje kapacitou svévolně zasáhnout.
 Rozdíl mezi oběma pojetími odhaluje klasický příklad, jenž stojí v jádru celé tradice: otrok laskavého pána. Pán, který svého otroka nikdy netrestá a ponechává mu široký prostor jednání, do jeho voleb fakticky nezasahuje — v liberálním smyslu je tedy otrok \uv{svobodný}.
  A přece cítíme, že svobodný není. Chybí mu nikoli nezasahování, nýbrž jistota: prostor, jejž požívá, není jeho právem, nýbrž pánovou libovůlí, kterou lze kdykoli odvolat.
   Otrok je nucen žít v ustavičné závislosti na cizí dobré vůli — a právě tato závislost, nikoli uskutečněný zásah, je oním bezprávím, jež republikánská tradice nazývá nadvládou.
Lovett tuto strukturu upřesňuje trojicí podmínek: nadvláda předpokládá (i) závislostní vztah, z něhož nelze bez značných nákladů vystoupit, (ii) nerovnováhu moci mezi stranami a (iii) možnost svévolného uplatnění této moci.
 Rozhodující je třetí podmínka a spolu s ní výklad pojmu svévole. 
 Zde se obě verze tradice mírně rozcházejí a je užitečné rozdíl přiznat.
  Podle procedurálního čtení je moc svévolná tehdy, není-li vázána sdílenými pravidly, jimž se podrobený může dovolat. 
  Podle Pettitova náročnějšího čtení je svévolná tehdy, není-li nucena sledovat zájmy a názory těch, jichž se týká — bez ohledu na to, zda se náhodou drží pravidel či nikoli.
   Pro naši analýzu je nosné právě Pettitovo pojetí, neboť činí zřejmým, proč benevolence nestačí: moc, která se chová vstřícně, avšak činí tak z vlastního rozhodnutí a nikoli z donucení respektovat zájmy podrobených, zůstává mocí nadvládnou.
Odtud plyne rozhodující obrat celé práce. Uživatel algoritmické platformy je vůči provozující korporaci v témže postavení jako otrok vůči laskavému pánovi.
 I když konkrétní aplikace algoritmu není namířena proti němu, ba i když je infrastruktura nastavena způsobem, jenž mu subjektivně prospívá, netěší se svobodě v republikánském smyslu. 
 Rozhoduje se totiž v prostředí, jehož parametry určuje nevolený soukromý subjekt disponující kapacitou toto prostředí kdykoli svévolně změnit — změnit, co uživatel uvidí, koho nalezne, čí hlas k němu dolehne — aniž by tato kapacita podléhala kontrole těch, kdo jsou jí vystaveni. 
 Nadvláda tedy nespočívá v uskutečněné škodě, nýbrž v samotné existenci takové nevázané, nekontrolované a demokraticky neodpovědné moci nad infrastrukturou, v níž se veřejná racionalita jedině může uskutečňovat.
Tímto vymezením získává práce tři výhody, jež strukturují celý následující výklad.
 

Za prvé, obvinění z nadvlády je robustní vůči empirické revizi: i kdyby se prokázalo, že algoritmická média nezpůsobují měřitelnou kognitivní škodu, bezpráví nadvlády by trvalo, neboť nezávisí na tom, zda a jak je moc uplatněna. 
 

Za druhé, přesouvá otázku z roviny terapeutické (jak chránit jednotlivce před závislostí a manipulací) do roviny ústavně-politické (jak moc rozdělit a učinit ji odpovědnou). 


Za třetí, vysvětluje strukturální selhání stávající regulace: individualisticky pojatá ochrana práv — souhlas, přístup k datům, právo na vysvětlení — nemůže nadvládu odstranit, neboť ta je vztahem, a nikoli událostí.
 Informovaný souhlas jednotlivce nadvládný vztah neruší, nýbrž jej stvrzuje; k jeho odstranění je třeba změnit rozložení moci samo. 
 K institucionálním důsledkům tohoto zjištění se práce vrací ve třetí kapitole.
\chapter{EKONOMIKA POZORNOSTI A TRANSFORMACE VEŘEJNÉ SFÉRY }
\section{Veřejná racionalita jako filozoficko-politický koncept}


\chapter{Formát PDF/A}

Трансформация понятия «мышление»: от человеческого к вычислительному За историей технических проектов искусственного разума стоит глубокая концептуальная трансформация: само понятие «мышления» последовательно сужалось и переопределялось так, чтобы машина могла в него вписаться.

В античности и в значительной мере в Средневековье мышление понималось как нечто несводимое к логическому выводу: оно включало нравственное суждение, эстетическое восприятие, мудрость как способность различать добро и зло в конкретной ситуации.
Aristotel'evskij nous был прежде всего способностью к созерцанию высших истин, а не механизмом дедукции.
Сенека в «Нравственных письмах» настаивал на том, что разум человека «причастен нравственным категориям, которые уходят гораздо дальше, чем просто логическое суждение — добро, зло, красота, уродство».
Именно поэтому «то, что действительно отличает человека от животного это способность к сложнейшей операции соединения логического и нравственного, нравственному выбору как таковому».
Это отличительная черта человеческого сознания: «выбирать не только рационально, но и нравственно».
Картезианский механицизм произвёл первое принципиальное сужение: мышление стало отождествляться с рациональным, логическим компонентом — тем, что потенциально поддаётся формализации.
Нравственное, эмоциональное, эстетическое было выведено за скобки или атрофировалось к «страстям», подлежащим управлению.
Лейбницевский проект «исчисления разума» сделал следующий шаг: мышление стало отождествляться с формальными операциями над символами.
Логика Буля, а затем математическая логика Фреге и Рассела завершили этот процесс на уровне формальных систем.
Тьюринг поставил принципиальный вопрос: если мы не можем отличить ответы машины от ответов человека в диалоге, есть ли у нас основания отрицать, что машина мыслит?
Этот вопрос был революционным — но именно потому, что он предполагал редукцию мышления к поведению, наблюдаемому извне.
Внутренний опыт, нравственная укоренённость, интенциональность — всё это было вынесено за скобки как ненаблюдаемое и потому нерелевантное.
Сёрл показал несостоятельность этой редукции: его «Китайская комната» демонстрирует, что система может производить правильные ответы, не понимая их — не имея никакого доступа к смыслу манипулируемых символов (Minds, Brains and Science, 1984).
Дрейфус добавил к этому аргумент воплощённости: даже если отвлечься от нравственного измерения, человеческое познание неотделимо от телесного опыта, от «знания-как» в отличие от «знания-что», — и никакая формальная система не воспроизведёт эту укоренённость (What Computers Can't Do, 1972).
Тем не менее прагматическая, редукционистская концепция мышления-как-вычисления оказалась исторически победоносной, но не потому, что она философски убедительнее, а потому, что она технически продуктивна и экономически выгодна.
Как показывает Петрунин, идея ИИ «опирается на господствующие в обществе интуиции» (Петрунин, б/г, с. 12) — и господствующей интуицией нашего времени оказалась интуиция технократическая: разум есть то, что можно измерить, формализовать и воспроизвести.
Всё остальное либо несущественно, либо будет формализовано позже.
Последствия этой победы выходят далеко за пределы технических дискуссий.
Если мышление — это вычисление, то алгоритм может принимать решения вместо человека: о кредитоспособности, о рисках, о виновности, о пригодности к должности.
Если интеллект измеряется поведенческим тестом, то система, проходящая тест, обладает достаточным интеллектом для управления.
Если нравственное измерение мышления нерелевантно, то вопрос об этической ответственности алгоритма не возникает.
Олоф-Орс и Смит показывают, что именно это концептуальное наследие делает возможным антропоморфный дизайн как инструмент надзорного капитализма: система, воспроизводящая поведение разумного существа, воспринимается пользователем как разумное существо и весь аппарат социального доверия, предназначенный для взаимодействия с другими людьми, оказывается задействован в интересах корпорации (Olof-Ors \& Smit, 2025).
В история от Луллия до ChatGPT мы наблюдаем последовательное переопределение мышления, чтобы машина могла в него вписаться.
На каждом этапе что-то важное исключалось из понятия: сначала нравственное суждение, потом интенциональность, потом воплощённость, наконец — само понимание.
Что осталось?
Статистически правдоподобная последовательность токенов.
Но это и не то, что философская традиция называла мышлением.
И это различие отнюдь не академическое: от него зависит, кто несёт ответственность за решения, принимаемые «интеллектуальными» системами, и чьи интересы эти решения обслуживают.

/chapter
\include{kap04}

\chapter*{Závěr}
\addcontentsline{toc}{chapter}{Závěr}

Předkládaná práce si položila otázku, jak algoritmicky zprostředkovaná digitální infrastruktura mění podmínky možnosti veřejné racionality a jaký druh politického problému tato změna představuje. 
Odpověď, k níž jsem dospěla, lze shrnout do čtyř tezí.


Za prvé: ztotožnění myšlení s výpočtem je pojmovým předpokladem algoritmické moci, nikoli jejím důsledkem. Rekonstrukce provedená v první kapitole ukázala, že cesta od Lullovy Ars Magna k současným jazykovým modelům je dějinami postupného vylučování: z pojmu myšlení byl nejprve vyloučen mravní úsudek, poté intencionalita, poté vtělenost a nakonec samotné porozumění. 
Politický význam tohoto vyloučení nespočívá v tom, že by stroje „nemyslely“ (spor o to nechávám otevřený) nýbrž v tom, že takto zúžený pojem rozumu učinil myslitelným delegování rozhodnutí na entitu, jež nesplňuje podmínky odpovědného subjektu.
 Depolitizace tedy není vedlejším účinkem technického pokroku; je zabudována do pojmové architektury, jež jej umožnila.


 Za druhé: platformový kapitalismus systematicky podkopává materiální předpoklady veřejné racionality, avšak rozsah této škody je sporný a nemůže nést tíhu politického argumentu. 
Druhá kapitola ukázala mechanismy: komodifikaci pozornosti, ohrazování zkušenosti do podoby dat, reálné podřízení komunikace logice kapitálu, erozi sdíleného epistemického prostoru.
 Zároveň jsem se však nemohla vyhnout tomu, že řada empirických tvrzení, o něž se kritická literatura opírá — kognitivní atrofie, epistemické bubliny, účinnost mikrotargetingu — je v odborné diskusi zpochybňována.
 Poctivá odpověď zní: nevíme s jistotou, jak velká ta škoda je. 
 Kritika, jež na těchto tvrzeních staví, je proto vratká.


 Za třetí, a to je hlavní teze práce: politické bezpráví algoritmické moci nespočívá ve škodě, nýbrž v dominanci.
Ať už je měřitelný kognitivní dopad platforem jakýkoli, trvá skutečnost, že nevolené soukromé subjekty disponují svévolnou, neprůhlednou a demokraticky nekontrolovanou mocí nad infrastrukturou, v níž se veřejná racionalita jedině může uskutečňovat: rozhodují, co je viditelné, komu jsou přiznány zdroje a příležitosti, a jaké epistemické prostředí občan obývá.
 V republikánském pojetí svobody jako nepřítomnosti dominance je toto bezpráví plné a nezávislé na tom, zda je moc uplatňována laskavě.
 Právě tento posun umožňuje vysvětlit, proč individualisticky pojatá ochrana práv — souhlas, přístup k datům, právo na vysvětlení — strukturálně selhává: dominance je vztah, nikoli událost, a nelze ji odstranit ani sebedokonalejším informovaným souhlasem jednotlivce.


 Za čtvrté: mezera odpovědnosti není technickým defektem, jenž by šlo technicky zaplnit.
 Rozostření odpovědnosti mezi vývojáře, provozovatele, zadavatele a regulátora reprodukuje logiku, kterou Arendtová popsala jako banalitu zla — nikoli jako historickou charakteristiku Eichmanna, nýbrž jako ideální typ byrokratické atomizace úkolu, v níž každý plní svou dílčí povinnost a souhrnný výsledek nemá původce.
 Antropomorfizace algoritmických systémů tuto mezeru nezaplňuje, nýbrž maskuje: připisuje subjektivitu artefaktu, jenž nesplňuje její podmínky, a osvobozuje tak skutečné subjekty od reflexe vlastního jednání. Odpovědnost však nelze delegovat na to, co nemá tvář.
\section*{Co z toho plyne}
Z takto vymezené diagnózy plynou tři odpovědi, jež se navzájem nevylučují, nýbrž doplňují. Ochrana (GDPR, AI Act) je nezbytná, avšak reaktivní, individualistická a asymetrická — a jak ukazuje osud tzv. 
Digital Omnibus, jímž bylo v roce 2026 uplatňování klíčových ustanovení AI Actu odloženo dříve, než vůbec vstoupilo v účinnost, je také vystavena systematickému tlaku regulovaných.
 Rozdělení moci — antimonopolní politika, povinný audit, oddělení infrastruktury od služeb — odpovídá povaze problému lépe, neboť míří na strukturu, nikoli na jednotlivý účinek.
 A konečně zmocnění: budování alternativní, federativní a společně vlastněné digitální infrastruktury, jež logiku dominance nereprodukuje.
Právě zde je však na místě střízlivost, k níž mne analýza dovedla. 
Existující alternativy — Fediverse, Wikipedia, otevřené jazykové modely, participativní platformy typu Decidim — dokládají, že jiná institucionální forma je možná. Nedokládají však, že je dostatečná.
 Migrace uživatelů na Mastodon po roce 2022 se z hlediska udržení uživatelů ukázala jako podstatně méně úspěšná, než naznačovala její první vlna; 
 participace na deliberativních platformách zůstává nízká; otevřené modely jsou závislé na výpočetní kapacitě, kterou nekontrolují. 
 Alternativní infrastruktura je nutnou, nikoli postačující podmínkou. 
 To není důvod k rezignaci, nýbrž k přesnějšímu zadání: otázka nezní, zda lze postavit lepší platformu, nýbrž jaké institucionální a majetkové uspořádání by zabránilo tomu, aby se moc znovu zkoncentrovala.
\section*{Meze práce a další výzkum}
Práce je prací pojmovou; empirické otázky, jichž se dotýká, nemůže rozhodnout.
 Přijímá habermasovský normativní rámec, ačkoli ten je sám předmětem sporu — agonistická kritika (Mouffe) by namítla, že ideál racionálního konsensu depolitizuje stejně účinně jako technokracie, proti níž je namířen; 
 tuto námitku jsem mohla pouze naznačit, nikoli vyřešit. 
Neanalyzuje mimoevropské konfigurace algoritmické moci ani otázku výpočetní suverenity, jež se pravděpodobně stane ústředním politickým problémem následujícího desetiletí.
Přesto se domnívám, že hlavní výsledek obstojí. Otázka, kterou algoritmická moc klade demokracii, není otázkou psychologickou ani technickou.
 Je to nejstarší otázka politické filozofie, položená v novém materiálu: kdo drží moc, komu se z ní zodpovídá a jak ji lze zpochybnit. 
 Kantovo Sapere aude nebylo adresováno géniovi schopnému samostatné myšlenky za jakýchkoli podmínek, nýbrž obyčejnému občanu, jehož schopnost k vlastnímu úsudku je strukturálně závislá na podmínkách, za nichž se formuje.
  Tyto podmínky dnes nevytváří veřejnost. Vytváří je několik soukromých korporací — a činí tak bez pověření, bez odpovědnosti a bez možnosti odvolání. 
  Obnovit veřejnou racionalitu proto neznamená naučit občany lépe zacházet s technologií. 
  Znamená to vrátit rozhodování o infrastruktuře veřejného rozumu do rukou těch, jichž se týká.


%%% Seznam použité literatury
\include{literatura}

%%% Obrázky v práci
%%% (pokud jich je malé množství, obvykle není třeba seznam uvádět)
\listoffigures

%%% Tabulky v práci (opět nemusí být nutné uvádět)
%%% U matematických prací může být lepší přemístit seznam tabulek na začátek práce.
\listoftables

%%% Použité zkratky v práci (opět nemusí být nutné uvádět)
%%% U matematických prací může být lepší přemístit seznam zkratek na začátek práce.
\chapwithtoc{Seznam použitých zkratek}

%%% Součástí doktorských prací musí být seznam vlastních publikací
\ifx\ThesisType\TypePhD \chapwithtoc{Seznam publikací} \fi

%%% Přílohy k práci, existují-li. Každá příloha musí být alespoň jednou
%%% odkazována z vlastního textu práce. Přílohy se číslují.
%%%
%%% Do tištěné verze se spíše hodí přílohy, které lze číst a prohlížet (dodatečné
%%% tabulky a grafy, různé textové doplňky, ukázky výstupů z počítačových programů,
%%% apod.). Do elektronické verze se hodí přílohy, které budou spíše používány
%%% v elektronické podobě než čteny (zdrojové kódy programů, datové soubory,
%%% interaktivní grafy apod.). Elektronické přílohy se nahrávají do SISu.
%%% Povolené formáty souborů specifikuje opatření rektora č. 72/2017.
%%% Výjimky schvaluje fakultní koordinátor pro zavěrečné práce.
\appendix
\chapter{Přílohy}

\section{První příloha}

\end{document}
